\documentclass[11pt]{article}
\usepackage[margin=1in]{geometry}
\usepackage{amsmath,amssymb,amsthm,mathtools}
\usepackage{enumitem}
\usepackage{hyperref}
\hypersetup{colorlinks=true,linkcolor=blue,citecolor=blue,urlcolor=blue}

\newtheorem{theorem}{Theorem}
\newtheorem{proposition}{Proposition}
\newtheorem{lemma}{Lemma}
\newtheorem{corollary}{Corollary}
\newtheorem{definition}{Definition}
\newtheorem{hypothesis}{Hypothesis}
\newtheorem{remark}{Remark}
\newtheorem{warning}{Warning}
\newtheorem{conjecture}{Conjecture}
\newtheorem{problem}{Problem}

\newcommand{\Z}{\mathbb Z}
\newcommand{\Q}{\mathbb Q}
\newcommand{\F}{\mathbb F}
\newcommand{\C}{\mathbb C}
\newcommand{\G}{\Gamma}
\newcommand{\E}{\mathbb E}
\newcommand{\1}{\mathbf 1}
\newcommand{\ord}{\operatorname{ord}}
\newcommand{\lcm}{\operatorname{lcm}}
\newcommand{\im}{\operatorname{im}}
\newcommand{\rank}{\operatorname{rank}}
\newcommand{\Gal}{\operatorname{Gal}}

\title{CRT Synchronization and Projective Kummer Slopes in Two-Generator Subgroups}
\author{Jared Wilder}
\date{Draft 0.4 --- 7-ROUND CONVERGED + R753 rank-$r$ section --- 2026-05-13}

\begin{document}
\maketitle

\begin{abstract}
For squarefree \(d\), membership in the subgroup \(\langle a,b\rangle\subset(\Z/d\Z)^*\) is not determined solely by membership in each prime component.  The same exponent vector must synchronize across all Chinese-remainder factors.  We formulate this obstruction as a lattice-image condition, define the synchronization defect, and show that the two-prime \(\ell\)-primary defect is exactly a rank defect.  In the motivating case \(a=2,b=3\), this rank defect is equivalent to a collision of projective Kummer slopes
\[
[\log_q2:\log_q3]\in\mathbb P^1(\F_\ell).
\]
\end{abstract}

\section{Introduction}

Let
\[
d=q_1\cdots q_r
\]
be squarefree.  A common mistake in two-generator congruence problems is to assume that local membership
\[
c\bmod q_i\in\langle a,b\rangle\subset\F_{q_i}^*
\]
for every \(i\) implies global membership
\[
c\bmod d\in\langle a,b\rangle\subset(\Z/d\Z)^*.
\]
This is false in general.  The same pair of exponents must work in every CRT component.

This paper records the algebraic obstruction.

\section{The lattice-image criterion}

Let \(d=q_1\cdots q_r\) be squarefree, and assume \(a,b,c\) are units modulo \(d\).  Choose primitive roots \(g_i\bmod q_i\) and write
\[
a=g_i^{u_i},\qquad b=g_i^{v_i},\qquad c=g_i^{w_i}.
\]

\begin{theorem}[Squarefree lattice membership]\label{T:lattice}
One has
\[
c\in\langle a,b\rangle\subset(\Z/d\Z)^*
\]
if and only if
\[
(w_1,\dots,w_r)\in
\im\left(
\Z^2\to\bigoplus_{i=1}^r\Z/(q_i-1)\Z
\right),
\]
where
\[
(k,\ell)\mapsto(ku_i+\ell v_i)_i.
\]
\end{theorem}

\begin{proof}
By CRT, the congruence \(a^k b^\ell\equiv c\pmod d\) is equivalent to the system
\[
a^k b^\ell\equiv c\pmod {q_i},\qquad i=1,\dots,r.
\]
In exponent coordinates this is
\[
ku_i+\ell v_i\equiv w_i\pmod {q_i-1}.
\]
The same pair \((k,\ell)\) must solve all congruences simultaneously, which is precisely the stated image condition.
\end{proof}

\section{Synchronization defect}

\begin{definition}[Synchronization defect quotient]
Let
\[
\Gamma_{q_i}(a,b)=\langle a,b\rangle\subset\F_{q_i}^*
\]
and use CRT to identify
\[
(\Z/d\Z)^*\simeq\prod_i\F_{q_i}^*.
\]
Let
\[
\Pi_d(a,b)=\prod_i\Gamma_{q_i}(a,b)
\]
and let
\[
\Gamma_d(a,b)=\langle a,b\rangle\subset \Pi_d(a,b)
\]
be the subgroup generated by the single global pair \(a,b\), embedded through CRT. Define the synchronization defect quotient by
\[
\mathcal S_d(a,b)=\Pi_d(a,b)/\Gamma_d(a,b).
\]
\end{definition}

\begin{proposition}[Synchronization size identity]
The synchronization defect quotient is finite and satisfies
\[
|\mathcal S_d(a,b)|
=
\frac{\prod_i|\Gamma_{q_i}(a,b)|}{|\Gamma_d(a,b)|}.
\]
Equivalently,
\[
|\Gamma_d(a,b)|=
\frac{\prod_i|\Gamma_{q_i}(a,b)|}{|\mathcal S_d(a,b)|}.
\]
\end{proposition}

\begin{proof}
By construction, \(\Gamma_d(a,b)\) is a subgroup of \(\Pi_d(a,b)\). The quotient is finite, and Lagrange's theorem gives the displayed formula.
\end{proof}

\begin{remark}
This quotient measures synchronization failure: \(\Pi_d(a,b)\) records independent local choices, while \(\Gamma_d(a,b)\) records choices coming from one common global exponent vector.
\end{remark}

\section{The two-prime rank criterion}

Now let \(d=pq\), and fix a prime \(\ell\mid\gcd(p-1,q-1)\).  Choose primitive roots and write
\[
a=g_p^{u_p}=g_q^{u_q},\qquad b=g_p^{v_p}=g_q^{v_q}.
\]
Modulo \(\ell\), form
\[
M_\ell(p,q)=
\begin{pmatrix}
u_p&v_p\\
u_q&v_q
\end{pmatrix}.
\]

\begin{theorem}[Pair rank defect]\label{T:rank}
The \(\ell\)-primary two-prime synchronization map has rank \(<2\) if and only if
\[
\det M_\ell(p,q)=u_pv_q-u_qv_p\equiv0\pmod\ell.
\]
\end{theorem}

\begin{proof}
This is the determinant criterion for a \(2\times2\) matrix over the field \(\F_\ell\).
\end{proof}

\section{Projective Kummer slopes}

Specialize now to \(a=2,b=3\).  For \(q\equiv1\pmod\ell\), choose a primitive root \(g_q\) and write
\[
2=g_q^{u_q},\qquad 3=g_q^{v_q}.
\]
Define the Kummer vector
\[
x_\ell(q)=(u_q,v_q)\in\F_\ell^2
\]
and, when nonzero, the projective Kummer slope
\[
\sigma_\ell(q)=[u_q:v_q]\in\mathbb P^1(\F_\ell).
\]

\begin{corollary}[Projective slope collision]\label{C:slope}
For primes \(p,q\equiv1\pmod\ell\), the rank defect in Theorem \ref{T:rank} occurs if and only if
\[
\sigma_\ell(p)=\sigma_\ell(q)
\]
in \(\mathbb P^1(\F_\ell)\), when both vectors are nonzero. The case where one or both vectors are \((0,0)\) is the star-atom case and is treated separately.
\end{corollary}

\begin{proof}
The determinant
\[
u_pv_q-u_qv_p
\]
vanishes exactly when the two nonzero vectors \((u_p,v_p)\) and \((u_q,v_q)\) span the same line in \(\F_\ell^2\).
\end{proof}

\section{Slope classes and Kummer ratio fields}

\begin{proposition}[Projective ratio field]\label{P:ratio}
For \([A:B]\in\mathbb P^1(\F_\ell)\), the condition
\[
[\log_q2:\log_q3]=[A:B]
\]
is equivalent, outside the star atom, to
\[
2^B3^{-A}\in(\F_q^*)^\ell.
\]
Thus slope classes are controlled by one-dimensional Kummer fields
\[
\Q(\zeta_\ell,(2^B3^{-A})^{1/\ell}).
\]
\end{proposition}

\begin{proof}
The equality of projective slopes is equivalent to
\[
B\log_q2-A\log_q3\equiv0\pmod\ell,
\]
which is equivalent to
\[
\log_q(2^B3^{-A})\equiv0\pmod\ell.
\]
That is exactly the condition that \(2^B3^{-A}\) is an \(\ell\)-th power modulo \(q\).
\end{proof}

\section{The star atom}

\begin{definition}
The star atom is the condition
\[
(\log_q2,\log_q3)=(0,0)\in\F_\ell^2.
\]
Equivalently, both \(2\) and \(3\) are \(\ell\)-th powers modulo \(q\).
\end{definition}

\begin{lemma}[Rational Kummer independence]\label{L:indep}
If
\[
2^a3^b\in(\Q(\zeta_\ell)^*)^\ell,\qquad 0\le a,b<\ell,
\]
then \(a=b=0\).
\end{lemma}

\begin{proof}
Suppose \(2^a3^b=\alpha^\ell\) in \(\Q(\zeta_\ell)\).  Take valuations at prime ideals above \(2\).  The right-hand side has valuations divisible by \(\ell\), so \(a\equiv0\pmod\ell\).  Since \(0\le a<\ell\), \(a=0\).  Similarly, valuations at prime ideals above \(3\) force \(b=0\).
\end{proof}

\begin{corollary}
The extension
\[
\Q(\zeta_\ell,2^{1/\ell},3^{1/\ell})/\Q(\zeta_\ell)
\]
has degree \(\ell^2\), and the expected relative Chebotarev density of the star atom among \(q\equiv1\pmod\ell\) is \(\ell^{-2}\).
\end{corollary}

\begin{proof}
The independence lemma shows that the classes of \(2\) and \(3\) span a two-dimensional subspace of \(K^*/(K^*)^\ell\), where \(K=\Q(\zeta_\ell)\).  Kummer theory gives relative degree \(\ell^2\).  The expected density statement is the Chebotarev density for complete splitting in this rank-two Kummer layer.
\end{proof}

\section{Collision graph}

\begin{definition}
For fixed \(\ell,Q\), define the Kummer slope collision graph \(G_\ell(Q)\) as follows.  Vertices are primes
\[
q\le Q,\qquad q\equiv1\pmod\ell.
\]
Two vertices are adjacent if their projective Kummer slopes agree, with the star atom counted separately.
\end{definition}

If \(N_s\) counts primes of slope \(s\in\mathbb P^1(\F_\ell)\), and \(N_\star\) counts the star atom, then
\[
|E(G_\ell(Q))|=
\sum_{s\in\mathbb P^1(\F_\ell)}\binom{N_s}{2}
+
\binom{N_\star}{2}.
\]



\section{Worked finite examples}

\subsection{Local membership without global synchronization}

Take \(d=5\cdot19=95\), \(a=2\), \(b=3\).  Use primitive root \(2\) modulo both \(5\) and \(19\).  Then
\[
2=2^1,\qquad 3=2^3\pmod5,
\]
so the exponent row modulo \(5\) is
\[
(u_5,v_5)=(1,3)\pmod4.
\]
Modulo \(19\),
\[
2=2^1,\qquad 3=2^{13}\pmod{19},
\]
so
\[
(u_{19},v_{19})=(1,13)\pmod{18}.
\]

Let \(c\) be the CRT class
\[
c\equiv1\pmod5,\qquad c\equiv2\pmod{19}.
\]
One representative is
\[
c=21\pmod{95}.
\]
Locally, \(c\) lies in \(\langle2,3\rangle\) modulo both \(5\) and \(19\).  In exponent coordinates, this asks for
\[
k+3\ell\equiv0\pmod4,
\]
and
\[
k+13\ell\equiv1\pmod{18}.
\]
Reducing the first congruence modulo \(2\) gives
\[
k+\ell\equiv0\pmod2,
\]
while reducing the second modulo \(2\) gives
\[
k+\ell\equiv1\pmod2.
\]
This is impossible.  Therefore \(c\) is locally represented at each prime factor but not globally represented modulo \(95\).

\subsection{Pair determinant zero}

Let \(p=7\), \(q=37\), and \(\ell=3\).  With primitive roots \(3\bmod7\) and \(2\bmod37\),
\[
(u_7,v_7)=(2,1)\pmod3,
\]
and
\[
(u_{37},v_{37})=(1,26)\equiv(1,2)\pmod3.
\]
The determinant is
\[
2\cdot26-1\cdot1=51\equiv0\pmod3.
\]
Thus the projective Kummer slopes collide modulo \(3\).

\subsection{Pair determinant nonzero}

Let \(p=7\), \(q=13\), and \(\ell=3\).  With primitive roots \(3\bmod7\) and \(2\bmod13\),
\[
(u_7,v_7)=(2,1)\pmod3,
\]
and
\[
(u_{13},v_{13})=(1,4)\equiv(1,1)\pmod3.
\]
The determinant is
\[
2\cdot4-1\cdot1=7\equiv1\pmod3,
\]
so the pair has no rank defect modulo \(3\).



\section{Verification transcript}

The included script \texttt{verify\_finite\_examples.py} verifies the displayed examples:
\[
\langle2,3\rangle\bmod23\text{ has size }11,
\]
\[
(u_5,v_5)=(1,3),\qquad (u_{19},v_{19})=(1,13),
\]
and the system
\[
k+3\ell\equiv0\pmod4,\qquad k+13\ell\equiv1\pmod{18}
\]
has no solutions. It also verifies
\[
\det M_3(7,37)\equiv0\pmod3
\]
and
\[
\det M_3(7,13)\not\equiv0\pmod3.
\]


\section{Additional examples to compute before submission}

A polished version of this paper should include explicit small examples:
\begin{enumerate}[leftmargin=1.5em]
\item a squarefree \(d=pq\) where local membership holds at \(p\) and \(q\) but global membership fails;
\item a pair \(p,q\) with \(\det M_\ell(p,q)=0\);
\item a pair \(p,q\) with \(\det M_\ell(p,q)\ne0\);
\item a small \(\ell\) example displaying the star atom.
\end{enumerate}
These are finite computations and do not require analytic number theory.

\section{Conclusion}

The synchronization obstruction is algebraic.  Analytic number theory enters only when one asks how often the projective Kummer slopes collide among primes.  That analytic question is separated from the algebraic content here.

\section{Rank-\(r\) synchronization}\label{S:rank-r}

The two-generator theory above extends to any number of generators.  Let \(a_1,\ldots,a_r\) be units modulo a squarefree integer \(d=q_1\cdots q_s\).  Choose primitive roots \(g_i\bmod q_i\), and write
\[
a_j\equiv g_i^{u_{ij}}\pmod{q_i},\qquad c\equiv g_i^{w_i}\pmod{q_i}.
\]
Define the rank-\(r\) exponent-lattice map
\[
\Phi_d:\Z^r\longrightarrow\bigoplus_{i=1}^s\Z/(q_i-1)\Z,
\qquad
(n_1,\ldots,n_r)\longmapsto
\left(\sum_{j=1}^r n_j u_{ij}\right)_{i=1}^s.
\]
Then
\[
c\in\langle a_1,\ldots,a_r\rangle\bmod d
\quad\Longleftrightarrow\quad
(w_i)_{i=1}^s\in\im(\Phi_d).
\]
This is the direct rank-\(r\) analogue of Theorem~\ref{T:lattice}; the proof is the same CRT-and-discrete-log argument applied to one global exponent vector \(n\in\Z^r\).

For \(q\equiv1\pmod\ell\), define the rank-\(r\) Kummer log row
\[
x_\ell(q)=(u_{q1},\ldots,u_{qr})\in\F_\ell^r.
\]
When \(x_\ell(q)\ne0\), the projective Kummer point is
\[
\sigma_\ell(q;a_1,\ldots,a_r)=[u_{q1}:\cdots:u_{qr}]\in\mathbb P^{r-1}(\F_\ell).
\]
For two primes \(p,q\equiv1\pmod\ell\) outside the star atom \(\{x_\ell(p)=0\}\cup\{x_\ell(q)=0\}\),
\[
\rank
\begin{pmatrix}
x_\ell(p)\\
x_\ell(q)
\end{pmatrix}<2
\quad\Longleftrightarrow\quad
\sigma_\ell(p)=\sigma_\ell(q)\ \text{in}\ \mathbb P^{r-1}(\F_\ell).
\]
The two-generator determinant criterion of Theorem~\ref{T:rank} is the \(r=2\) case of projective collision in \(\mathbb P^{r-1}(\F_\ell)\).

The companion paper \emph{Rank-\(r\) synchronization defects and projective Kummer geometry} (Wilder, 2026) develops the full rank-\(r\) calculus, including the synchronization defect quotient size formula \(|\Gamma_d|=\prod_i|\Gamma_{q_i}|/|\mathcal S_d|\), the first-level \(\ell\)-defect cokernel and its rank-defect lower bound \(\delta_\ell(T;\mathbf a)\ge |T|-r\), the hyperplane atom criterion \(A\cdot x_\ell(q)=0\iff\alpha_A\in(\F_q^*)^\ell\) with controlling field \(K_{\ell,A}=\Q(\zeta_\ell,\alpha_A^{1/\ell})\), the rank-\(r\) star field of degree \(\ell^r\) with expected density \(\ell^{-r}\), and the rank-defect hypergraph \(H^{(k)}_{\ell,r}(Q;\mathbf a)\).

\bibliographystyle{alpha}
\begin{thebibliography}{9}

\bibitem{Lang}
S. Lang.
\newblock \emph{Algebra}.
\newblock Springer.

\bibitem{Washington}
L. Washington.
\newblock \emph{Introduction to Cyclotomic Fields}.
\newblock Springer.

\end{thebibliography}

\end{document}
